Este capítulo é reservado à análise dos dados do instrumento descrito na Seção~\ref{sec:protocolo-avaliacao}. Até o fechamento desta revisão a coleta não tinha sido feita: o formulário está escrito e revisado, mas não foi distribuído, e o encaminhamento ao comitê de ética da instituição continua pendente. Por isso não há resultados de participantes aqui.

As seções abaixo definem desde já o que será relatado e com quais critérios. Se essa definição fosse feita depois de ver os dados, seria possível escolher o recorte que deixasse a plataforma em melhor situação; defini-la agora evita essa escolha. Os critérios de cálculo dos escores estão na Seção~\ref{sec:protocolo-avaliacao} e não são repetidos aqui.

% =====================================================================
% A PREENCHER APÓS A COLETA.
% Cada seção abaixo diz o que relatar e com qual critério. Manter a ordem
% e os critérios definidos aqui; se algo mudar depois de ver os dados, a
% mudança precisa ser declarada no texto.
% =====================================================================

\section{Composição do Painel}
\label{sec:resultado-painel}

% A PREENCHER: quantos responderam; distribuição por titulação, área de
% atuação e tempo de experiência; quantos lecionam ou já lecionaram; quais
% ferramentas de ensino o painel já conhecia.
% Numa avaliação feita por especialistas, o perfil de quem avalia faz parte
% do resultado: apresentar antes dos escores.
% Registrar também o que os respondentes disseram ter feito na plataforma e
% o que personalizaram. Sem isso não dá para saber se uma avaliação
% negativa veio de quem usou o assistente inteiro ou de quem nem o abriu.

\section{Assistente de Personalização da Linguagem}
\label{sec:resultado-assistente}

% A PREENCHER: itens do bloco de recursos sobre o assistente -- se foi
% simples criar ou ajustar uma linguagem, se ficou claro em que etapa o
% usuário estava e se a pré-visualização ajudou.
% Comparar cada escore com a decisão de interface correspondente descrita
% na Seção \ref{sec:customizacao-dinamica}: esses itens verificam escolhas
% adotadas como hipótese, e um resultado ruim indica qual decisão revisar.
% Relatar à parte quantos marcaram "não usei" em cada linha: isso mede
% quantas pessoas chegaram ao recurso, e não a qualidade dele.

\section{Uso do Ambiente Virtual de Aprendizagem}
\label{sec:resultado-ava}

% A PREENCHER: bloco respondido por quem leciona ou já lecionou -- criar
% turma, criar exercício com casos de teste, travar linguagem no exercício,
% correção automática, preparar aula sem apoio técnico e relação entre o
% esforço de preparação e o ganho em sala.
% Incluir as respostas sobre o momento da disciplina em que a plataforma
% seria mais útil, sem esconder a alternativa que diz não haver espaço para
% ela: se muitos marcarem essa opção, isso é um resultado, não um erro.

\section{Usabilidade e Experiência de Uso}
\label{sec:resultado-usabilidade}

% A PREENCHER: escore SUS geral, com intervalo de confiança só se houver
% vinte respondentes ou mais; com menos, apresentar como indicativo e não
% afirmar que há diferença entre subgrupos. Comparar com a média de
% referência e com as faixas de interpretação adotadas.
% UEQ-S: relatar as dimensões pragmática e hedônica separadas, na faixa de
% -3 a +3. Se elas divergirem, discutir: uma plataforma considerada
% eficiente mas pouco interessante (ou o contrário) diz algo sobre a
% interface que o escore SUS sozinho não mostra.

\section{Potencial Pedagógico e Intenção de Uso}
\label{sec:resultado-pedagogico}

% A PREENCHER: bloco de potencial pedagógico (oito itens em terceira
% pessoa) e bloco vindo do TAM.
% As três primeiras afirmações do bloco pedagógico correspondem às
% premissas da introdução. Verificar se o painel concorda com elas e dizer
% com clareza caso não concorde.

\section{Riscos e Objeções Levantados pelo Painel}
\label{sec:resultado-riscos}

% A PREENCHER: bloco de riscos, relatado SEPARADAMENTE e avisando o leitor
% de que ali concordar significa apontar um problema.
% Não somar nem tirar média junto com o bloco de potencial pedagógico: as
% escalas apontam em sentidos opostos e uma anularia a outra.
% Para cada risco apontado, dizer se a implementação atual já resolve (o
% travamento de linguagem por exercício e por lista, por exemplo, responde
% à objeção sobre turmas com sintaxes diferentes) ou se continua em aberto.

\section{Respostas Abertas}
\label{sec:resultado-abertas}

% A PREENCHER: análise das cinco questões abertas -- o que funcionou
% melhor, onde a pessoa teve dificuldade ou se perdeu, quais mensagens de
% erro foram confusas e o que mudaria na plataforma.
% Citar trechos apenas de quem autorizou a citação, sem identificação,
% conforme o item de autorização do formulário.
% Os relatos de quem se perdeu interessam bastante a este trabalho: eles
% mostram onde a navegação da interface falhou.

\section{Síntese e Confronto com os Objetivos}
\label{sec:resultado-sintese}

% A PREENCHER: retomar os objetivos específicos e dizer, para cada um, o
% que a avaliação sustenta e o que continua sem verificação.
% Manter a distinção da Subseção \ref{subsec:limites-interpretacao}: o que
% foi coletado é a opinião de profissionais sobre a plausibilidade da
% redução da barreira sintática, e não a prova dessa redução. Afirmar mais
% do que isso vai além do que o estudo permite.
